% =====================================================================
% Section 01 — Introduction
% CONTENT PLAN: (1) problem: "scale it" conflates capabilities and is expensive
%   on two permanent axes (OpEx, CapEx). (2) Two deployment realities motivate
%   the map: regulated/closed -> small on-prem + compliance first-class (precise
%   regulatory framing: regs demand traceability/validation/oversight, NOT
%   determinism per se); per-domain authoring is the recurring cost -> fixed
%   scaffold + cheap A2 -> A2-swap. (3) the decomposition. (4) objective
%   function. (5) contributions (headline = capability vs compliance scaling).
% SUBSECTION TITLES: currently uses \paragraph runs, not \subsection. If the PI
%   wants numbered subsections (e.g. 1.1 Motivation, 1.2 Objective,
%   1.3 Contributions), set the titles here.
% =====================================================================
\section{Introduction}\label{sec:intro}
A tool-use agent maps a natural-language request to a sequence of tool calls that satisfy a domain policy. The default recipe for improving such an agent is a larger base model. But ``larger'' conflates distinct capabilities, and it is expensive on two permanent axes: inference cost (a model twice the size costs more than twice as much per request and is slower, for the lifetime of the deployment) and, for on-premises or sovereign deployment, hardware capital expenditure, since the operator must buy the memory, power, and cooling. The practical question is therefore not \emph{how large a model} but \emph{which capabilities must scale, which can be delegated to cheap deterministic structure, which no model can be trusted to perform, and what mixture minimizes total cost}.

\paragraph{Motivation: two deployment realities.}% TODO(아버지): subsection 제목 확정 가능
First, regulated and closed deployments force small on-premises models and make compliance a first-class concern. In regulated or sovereign settings---finance, healthcare, government---customer data cannot leave the premises, so a frontier API is unavailable and the operator runs, and buys the hardware for, a local model. At the same time, regulation increasingly demands traceability, validation, and effective human oversight (for example EU AI Act Articles 12 and 14, and US banking guidance SR~26-2)---but not determinism as such; no major framework textually mandates it, the medical-device provision EU MDR Annex~I~\S17.1 being the lone exception. A deterministic policy gate satisfies these obligations by construction, and under SR~26-2 deterministic rule-based processes fall outside the model-risk perimeter entirely, whereas a stochastic agent can demonstrate compliance only statistically, per sample, with its effective policy unverifiable in advance. Our central finding---that compliance is scale-invariant (\S\ref{sec:results-compliance})---implies that such a gate is needed at every scale, frontier models included.

Second, per-domain authoring is the recurring cost, so a fixed scaffold combined with cheap-to-author relations is the path to many domains. Domain-specific learning is slow and expensive and must be redone for each field; if the per-domain artifact (the A2) can be authored cheaply and the scaffold and learned skill are held fixed, then moving to a new domain costs an A2-swap rather than a retrain. This paper measures single-domain compliance and cost; the multi-domain payoff is the generalization axis of the cost framework (\S\ref{sec:knee}), which we validate only partially here (a SOPBench cross-domain check) and otherwise treat as a design premise rather than a measured multi-benchmark result. Taken together, a deployer must use a small on-premises model for cost and sovereignty, needs a form of compliance that scale does not provide, and wants to amortize effort across domains.

We answer with a capability $\times$ scale $\times$ lever $\times$ cost decomposition. For each functional sub-capability we ask, across scale, whether it (a)~grows with scale, (b)~is a genuine model limit, (c)~is covered by deterministic \emph{scaffold}, or (d)~is covered by per-domain authored relations (\emph{A2}); we then assemble the per-capability levers under a total-lifecycle cost objective and locate the cost knee. The benchmarks are instruments for filling these cells, not the object of study.

\paragraph{Objective.}% TODO(아버지): subsection 제목 확정 가능
We seek to minimize, simultaneously, the cost of cheap GPU tier and count, the memory footprint (favoring int8 or int4 quantization), the maintenance burden per change, the number of human experts required, and the cost of generalizing to each new field (ideally an A2-swap alone), together with hardware capital and inference operating expenditure. These terms are coupled and point in the same direction: cheap, low-memory hardware favors a small, quantized model, while low maintenance and onboarding favor maximizing the fixed domain-general component (deterministic scaffold plus transferable, forgetting-free learning) and minimizing the per-domain A2. The A2 is the common adversary, since it drives maintenance, experts, and per-field cost, and is worst on-premises where a frontier A2-generator is unavailable; a fixed scaffold and transferable skill are the common ally. The optimum is the total-cost knee, not the smallest model.

\paragraph{Contributions.}% TODO(아버지): subsection 제목 확정 가능
(1)~A controlled, frontier-API-free methodology that isolates operand execution, orchestration load, and compliance across scale (\S\ref{sec:method}).
(2)~Evidence that operand execution is scale-saturated and not a learnable gap at 32B (\S\ref{sec:results-operand}); the apparent gap is criterion interpretation and join/disambiguation. Faithful-formalization fine-tuning is therefore unwarranted at 32B, though a residual may exist at 7B (\S\ref{sec:future}).
(3)~A load decomposition with a scale response (\S\ref{sec:results-load}); controlled generation overturns a correlation visible only in observational data.
(4)~The finding that compliance is scale-invariant (\S\ref{sec:results-compliance})---the empirical backbone of a two-kind scaffold (a capability scaffold that recedes with scale, and a guarantee scaffold that does not).
(5)~A total-cost-optimal lever allocation with a deployment payoff (\S\ref{sec:results-cost}--\ref{sec:knee}): decidability-first routing, an on-premises cost-per-request advantage of roughly an order of magnitude (the exact multiple reported as an estimate pending measured token accounting), and a projected fleet design, framed by the cost knee.
(6)~A proposed unifying operator---deterministic scaffold as load reduction (\S\ref{sec:framework})---for which we establish the components but defer the full predictive test (\S\ref{sec:future}), and which we therefore present as a framework hypothesis rather than a verified law.
Our claim is deliberately \emph{not} ``small beats large,'' a result already established by reinforcement-learned routing systems (\S\ref{sec:related}); it is \emph{which} parts of the gap are scale-invariant and must be offloaded, and a cost-optimal map indicating which lever a deployer should pull for each capability.
